\begin{abstract}
% Rascunho inicial do resumo (professor). Reescrever com os n\'umeros
% finais do Exemplo 2 quando as simula\c{c}\~oes estiverem conclu\'idas.
A corros\~ao induzida por carbonata\c{c}\~ao \'e um dos principais mecanismos de
degrada\c{c}\~ao de estruturas de concreto armado, e a previs\~ao acurada da vida
\'util remanescente (\textit{Remaining Useful Life} --- RUL) \'e essencial para a
gest\~ao de infraestruturas informada por risco. Este trabalho prop\~oe um emulador
estoc\'astico dependente do tempo que aproxima a distribui\c{c}\~ao condicional
completa da fun\c{c}\~ao de estado limite $g(\mathbf{X},t)$ por meio da
Distribui\c{c}\~ao Lambda Generalizada, cujos par\^ametros s\~ao representados por
Expans\~ao em Caos Polinomial das vari\'aveis de entrada e interpolados
continuamente no tempo por um modelo de aprendizado de m\'aquina. O arcabou\c{c}o
integra o modelo de avan\c{c}o da frente de carbonata\c{c}\~ao, a perda de se\c{c}\~ao
e de ader\^encia das armaduras e a verifica\c{c}\~ao do estado limite, fornecendo,
a custo computacional desprez\'ivel, a probabilidade de falha, a distribui\c{c}\~ao
do tempo at\'e a falha, a RUL probabil\'istica e m\'etricas de risco de cauda (VaR e
CVaR). O m\'etodo \'e verificado em um problema de refer\^encia com solu\c{c}\~ao de
Monte Carlo direta como \textit{ground truth} e aplicado \`a previs\~ao da RUL de
elementos de concreto armado sob carbonata\c{c}\~ao. Os resultados indicam
excelente concord\^ancia distribucional (erro relativo inferior a
\falta{n\'umero}\% at\'e o percentil 99) e ganho de efici\^encia computacional de
\falta{n\'umero} ordens de grandeza em rela\c{c}\~ao \`a simula\c{c}\~ao direta,
viabilizando an\'alises de RUL e de cen\'arios de manuten\c{c}\~ao em tempo quase
real.
\end{abstract}

\textbf{Palavras-chave}: vida \'util remanescente; carbonata\c{c}\~ao; emulador
estoc\'astico; distribui\c{c}\~ao lambda generalizada; confiabilidade dependente do
tempo; concreto armado.

\vspace{0.5em}
\noindent\textbf{Keywords}: remaining useful life; carbonation; stochastic emulator;
generalized lambda distribution; time-dependent reliability; reinforced concrete.
